\section{The Frankfurt School}

It seems, judging from the among of attention it gets, that the Frankfurt school is better known, at least in name, among those on the ``right'' than those on the left who are allegedly influenced by its doctrines. Its founders were Horkheimer and Adorno, two basically unreadable authors. Then there was Herbert Marcuse, the idol of the New Left, and Eric Fromm, who was able to transpose many of their ideas into the idiom of popular culture.

The fundamental technique of the school is to combine the doctrines of Sigmund Freud and Karl Marx in a coherent system. It is important to understand the attraction of this doctrine. In any spiritual teaching --- and I include the Frankfurt school, since it reflects a battle of minds --- first the problem must be defined, and then a solution proposed. The problem is defined as widespread human misery resulting from inequalities. The solution, in conformity with the socialist mindset\footnote{See Section \ref{sec:LiberalismSocialism} in this book.}, is to locate it in social structures.

By means of Marxist theory, they move Freudianism out of the individual into the collectivity: the cause of oppression is not one's own superego, but rather the patriarchy that forms the superego. Thus, revolution against the established system of the patriarchy will lead to freedom; that is, the elimination of the internalized patriarchal values that create conflict in the individual, will lead to the free self-expression of the id in the pursuit of its instinctual pleasures. This revolutionary attitude is now commonplace, and even a ``moderate'' like Al Gore, can tell students to ignore their parents.

Obviously, this point of view is completely opposed to Tradition. In it, there is no transcendence and the primacy of instinct is the opposite of the control of the appetites by the spirit. Nevertheless, although the school is false in what it denies, there may be some truth in what it asserts. Even Julius Evola concedes that Freudian and Marxian analyses can apply to contemporary mass man as he is. Since the mass man has no sense of transcendence, particularly self-transcendence, nor any awareness of the role of Will in creating the circumstances of his life, he is open to see the sources to his basic dissatisfaction as coming from outside himself. The reinforcement of this attitude leads to resentment against hierarchical authority; the promise of the satisfaction of instinctual needs leads to support.

So, it should be obvious that the Frankfurt school in itself is not the problem. It is merely the latest phase of revolutionary thought that began with the Reformation and continued with the French and Russian revolutions. Since there are precious few that hold a complete counter-revolutionary mindset, there is little of intellectual depth able to point out where the school is false. There are attempts to oppose it on the basis of such factors as economic policies, genetics, religion (predominantly non-traditional), or race. However, none of these are ``radical'', that is, they do not get to the roots of proper thinking. Take this quiz\footnote{See Section \ref{sec:Counterrevolutionary} in this book.} again and see where you really stand.


\flrightit{Posted on 2010-04-01 by Cologero}

